\section{The manifold and the obstruction to nonpositive curvature}
\label{sec:topology}

We first choose the underlying manifold.  Its two product pieces impose
incompatible orthogonality conditions on the gluing torus under any
nonpositively curved metric.  This is an instance of Leeb's two-piece graph-manifold
obstruction \cite[Section~4.1, Example~4.1]{Leeb1995}.  The obstruction to
proper cocompact actions on proper CAT(0) spaces is also contained in the
group-theoretic framework of Kapovich and Leeb
\cite[Theorem~2.4]{KapovichLeeb1996}.  We give the group-theoretic argument
explicitly, since it fixes the gluing convention used in the metric construction.

\subsection{Two product pieces and one shear}

Let $\Sigma$ be a compact connected oriented surface of genus one with one
boundary component, and let $N_i=\Sigma\times S^1$ for $i=1,2$.
On $\partial N_i$ write $h_i$ for the oriented boundary circle of $\Sigma$
and $f_i$ for the positively oriented circle factor.  Thus $(h_i,f_i)$ is
an oriented integral basis of $\pi_1(\partial N_i)\cong\mathbb Z^2$.
Choose the linear diffeomorphism $\varphi:\partial N_2\longrightarrow
\partial N_1$ whose action on these bases is
\begin{equation}\label{eq:torus-gluing}
 h_2=h_1+f_1,\qquad f_2=h_1.
\end{equation}
Here and below additive notation is used only inside the torus group.
The matrix of $\varphi_*$ is
$\left(\begin{smallmatrix}1&1\\1&0\end{smallmatrix}\right)$, with determinant
$-1$, so $\varphi$ reverses the induced boundary orientations.
Gluing the two pieces by $\varphi$ therefore gives a closed connected
oriented smooth three-manifold $M$.  More explicitly, collar coordinates
give smooth charts across the identified boundary; the compact pieces
cover $M$, every boundary component is paired, and the common torus makes
the union connected.  We insert a product cylinder $[0,1]\times T^2$
between the two pieces.  Inserting this collar leaves the diffeomorphism
type unchanged and provides the region where the metric will vary.

The coordinate realization needed later is worth recording.  For any
$\epsilon>0$ and $B_2>0$, put $B_1=\epsilon B_2$ and
\[
 T^2=\mathbb R^2/\Lambda,\qquad
 \Lambda=\mathbb Z B_1e_1+\mathbb Z B_2e_2,
\]
where $e_1,e_2$ are the standard coordinate vectors.  At the first end use
the frame $(e_1,e_2)$ with respective coordinate periods $(B_1,B_2)$.
At the second end use the frame $(e_2+\epsilon e_1,e_1)$ with respective
periods $(B_2,B_1)$.  The first vector in each frame is the surface-boundary
direction and the second is the fiber direction.  Their period vectors are
\begin{equation}\label{eq:torus-periods}
 \begin{aligned}
 h_1&: B_1e_1,& f_1&:B_2e_2,\\
 h_2&:B_2(e_2+\epsilon e_1)=B_1e_1+B_2e_2,
 &f_2&:B_1e_1.
 \end{aligned}
\end{equation}
In particular the second pair is a basis of exactly the same lattice,
and these physical periods implement \eqref{eq:torus-gluing}.

Choose basepoints and connecting paths through the cylinder so that the
boundary groups are identified by \eqref{eq:torus-gluing}.  For suitable
free generators $a_i,b_i$ of $\pi_1(\Sigma)$, the boundary element is
$h_i=[a_i,b_i]^{\pm1}$.  This cyclically reduced nonempty word has infinite
order in the free group.  Consequently
\[
 G_i:=\pi_1(N_i)=F(a_i,b_i)\times\langle f_i\rangle,
 \qquad
 \langle h_i,f_i\rangle\cong\mathbb Z^2
\]
and the boundary-group inclusion in each $G_i$ is injective: a relation
$h_i^m f_i^n=1$ first gives $n=0$ by projection onto the circle factor
and then $m=0$ in the free group.  Van Kampen's theorem and the normal form
for an amalgamated product give
\begin{equation}\label{eq:amalgam}
 \Gamma:=\pi_1(M)=G_1*_{\Gamma_T}G_2,
 \qquad \Gamma_T=\langle h_1,f_1\rangle=\langle h_2,f_2\rangle\cong\mathbb Z^2,
\end{equation}
with both vertex groups and $\Gamma_T$ embedded in $\Gamma$.
For $f\in\Gamma$, write
$C_\Gamma(f)=\{g\in\Gamma:gf=fg\}$ for its centralizer.  Since $f_i$ is
central in $G_i$, we have
\begin{equation}\label{eq:boundary-commutators}
 G_i\subset C_\Gamma(f_i),\qquad
 h_i\in[C_\Gamma(f_i),C_\Gamma(f_i)]\quad(i=1,2).
\end{equation}

\subsection{The two incompatible orthogonality conditions}

\begin{theorem}\label{thm:no-cat0-action}
The fundamental group $\Gamma=\pi_1(M)$ of the closed three-manifold
defined by \eqref{eq:torus-gluing} admits no proper cocompact isometric
action on a nonempty proper $\operatorname{CAT}(0)$ space.
\end{theorem}

\begin{proof}
Suppose that $\Gamma$ acts properly and cocompactly by isometries on a
nonempty proper $\operatorname{CAT}(0)$ space $X$.  Properness of the action
means that for compact sets $K,L\subset X$, only finitely many
$g\in\Gamma$ satisfy $gK\cap L\ne\varnothing$; in particular, point
stabilizers are finite.

For $g\in\Gamma$ set
\[
 \ell(g)=\inf_{x\in X}d(x,gx),\qquad
 \operatorname{Min}(g)=\{x\in X:d(x,gx)=\ell(g)\}.
\]
We check the attainment of this infimum before applying the flat torus
theorem.  Choose $x_n$ with $d(x_n,gx_n)\to\ell(g)$ and a compact set
$K\subset X$ whose $\Gamma$-translates cover $X$.  There are $k_n\in\Gamma$
such that $y_n=k_nx_n\in K$.  Put $g_n=k_ngk_n^{-1}$.
For all large $n$, $d(y_n,g_ny_n)\le\ell(g)+1$.  The closed
$(\ell(g)+1)$-neighborhood of $K$ is compact, because $X$ is proper.
Properness of the action therefore permits only finitely many $g_n$.
After passage to a subsequence, $g_n=g_*$ and $y_n\to y\in K$, whence
$d(y,g_*y)=\ell(g)$.  Since translation length is invariant under
conjugation, conjugating $y$ back gives a point of
$\operatorname{Min}(g)$.  Thus every element of $\Gamma$ is semisimple.
If $g$ has infinite order, its attained minimum is positive: otherwise
$g$ fixes a point, and its infinitely many distinct powers belong to that
point stabilizer, contrary to its finiteness.  The semisimplicity argument
is also given in
\cite[Proposition~II.6.10(2)]{BridsonHaefliger1999}.

Apply the flat torus theorem
\cite[Theorem~II.7.1(1)--(3)]{BridsonHaefliger1999} to the embedded group
$\Gamma_T\cong\mathbb Z^2$.  Its action is proper and semisimple by the preceding
paragraph.  Hence
\[
 \operatorname{Min}(\Gamma_T):=\bigcap_{a\in\Gamma_T}\operatorname{Min}(a)
       \cong Y\times\mathbb R^2,
\]
and $\Gamma_T$ acts trivially on $Y$ and as a translation lattice on
$\mathbb R^2$.  Fix one of these invariant Euclidean planes, denoted $E$,
and write $\mathbf v(a)\in\mathbb R^2$ for the translation vector of
$a\in\Gamma_T$.  The map $\mathbf v$ is a homomorphism, and its values on any
integral basis of $\Gamma_T$ are linearly independent.

We now use the boundary-commutator relation
\eqref{eq:boundary-commutators}.  For each $i$, the positive translation
length $\ell_i=\ell(f_i)$ and the cyclic minset theorem give
\begin{equation}\label{eq:cyclic-minset}
 \operatorname{Min}(f_i)\cong Y_i\times\mathbb R,
 \qquad f_i(y,t)=(y,t+\ell_i).
\end{equation}
Every $g\in C_\Gamma(f_i)$ preserves this product and acts as
\[
 g(y,t)=(g_i'(y),t+\tau_i(g)),
\]
where $\tau_i(g)$ is a real number independent of $y$;
see \cite[Theorem~II.6.8(4)--(5)]{BridsonHaefliger1999}.
To see the reason for this product action, commutation sends each oriented
$f_i$-axis to another such axis.  In the splitting
\eqref{eq:cyclic-minset}, these are the lines $\{y\}\times\mathbb R$.
An axiswise translation could initially have a shift depending on $y$,
but preservation of the product distance between $(y,t)$ and $(z,s)$
for every $t-s\in\mathbb R$ forces the two shifts to agree.
Composition now shows that
$\tau_i:C_\Gamma(f_i)\longrightarrow(\mathbb R,+)$ is a homomorphism.
It vanishes on commutators, so $\tau_i(h_i)=0$.

The same plane $E$ lies in both cyclic minsets, because
$E\subset\operatorname{Min}(\Gamma_T)$.  At a point $x=(y,t)\in E$ expressed
in the $i$th splitting, the zero shift of $h_i$ gives
\[
 \begin{aligned}
 d(x,h_if_ix)^2
 &=d_{Y_i}(y,h_i'y)^2+\ell_i^2\\
 &=d(x,h_ix)^2+d(x,f_ix)^2.
 \end{aligned}
\]
All three distances can also be computed in the Euclidean plane $E$.
It follows that
\begin{equation}\label{eq:torus-orthogonality}
 \langle\mathbf v(h_i),\mathbf v(f_i)\rangle=0
 \qquad(i=1,2).
\end{equation}
Put $u=\mathbf v(h_1)$ and $v=\mathbf v(f_1)$.
The first equality in \eqref{eq:torus-orthogonality} says
$\langle u,v\rangle=0$.  The gluing relations give
$\mathbf v(h_2)=u+v$ and $\mathbf v(f_2)=u$, so the second equality says
\[
 0=\langle u+v,u\rangle=|u|^2.
\]
This is impossible: $h_1$ is nontrivial in $\Gamma_T$, and its translation vector
in the lattice is nonzero.
\end{proof}

\begin{corollary}\label{cor:no-local-cat0}
The manifold $M$ admits no locally $\operatorname{CAT}(0)$ length metric
inducing its manifold topology.
\end{corollary}

\begin{proof}
Suppose that such a metric exists.  Compactness of $M$ makes it complete
and locally compact.  Its ordinary universal cover, equipped with the
lifted length metric, is complete and locally compact, and the covering
map is a local isometry.  Metric Cartan--Hadamard makes this universal
cover a $\operatorname{CAT}(0)$ space, and metric Hopf--Rinow makes it
proper; see \cite[Proposition~I.3.7(1) and Theorem~II.4.1(2)]{BridsonHaefliger1999}.
The deck group $\Gamma$ acts freely and properly by isometries, with
compact quotient $M$.  This contradicts Theorem~\ref{thm:no-cat0-action}.
\end{proof}

\begin{proposition}\label{prop:no-npc}
The closed smooth three-manifold $M$ defined by
\eqref{eq:torus-gluing} admits no smooth Riemannian metric of nonpositive
sectional curvature.
\end{proposition}

\begin{proof}
Such a Riemannian metric is a length metric inducing the manifold topology
and is locally $\operatorname{CAT}(0)$
\cite[Theorem~II.1A.6]{BridsonHaefliger1999}, contradicting
Corollary~\ref{cor:no-local-cat0}.
\end{proof}

These obstructions depend only on the topology of $M$.  The metric
constructed below will have no conjugate points.
